% ch02_u05: 原书 p.53--p.60
% 首句承接 ch02_u04 末尾（原书 p.52 页末断句，同一段未完）。
\SourcePage{53}
代数群，作用在 $M$ 上。则 $T^{*}(H\backslash M)=H\backslash N$，其中
$N\subset T^{*}M$ 是 $H$ 各轨道的余法丛之并。每个余法丛的维数纯为
$\dim M$；又因 $g=0$，$H$-轨道的个数有限。

\emph{注.} 本质上相同的论证表明：对任何光滑代数栈 $\mathcal{Y}$，
$T^{*}\mathcal{Y}$ 在每一点的维数都 $\ge\dim\mathcal{Y}$。若
$g=0$ 且 $\mathcal{Y}=\operatorname{Bun}_G$，则按定理 2.10.2 有
$T^{*}\mathcal{Y}=\operatorname{Nilp}$，且
$\dim T^{*}\mathcal{Y}=\dim\mathcal{Y}$。于是又证得 $g=0$ 时的
定理 2.10.4。

2) 设 $g=1$。此时宜设 $G$ 约化，而不必半单（这其实并非本质：
约化 $G$ 的定理 2.10.4 易由半单情形推出）。

在进入证明之前，先回忆半稳定性与 Shatz 分层的概念。固定 Borel 子群
$B\subset G$，记 $H$ 为其极大交换商。设 $P\subset G$ 是含 $B$ 的
抛物子群，$L$ 为 $P$ 的极大约化商，$Z$ 为 $L$ 的中心。设 $\Gamma$
（resp. $\Delta$）为 $G$（resp. $L$）的单根集。嵌入
$Z\hookrightarrow L$ 诱导同构
\[
  \operatorname{Hom}(Z,\mathbb{G}_m)\otimes\mathbb{Q}
  \xrightarrow{\sim}
  \operatorname{Hom}(L,\mathbb{G}_m)\otimes\mathbb{Q}.
\]
记 $p$ 为合成
\[
  \operatorname{Hom}(H,\mathbb{G}_m)\to
  \operatorname{Hom}(Z,\mathbb{G}_m)\to
  \operatorname{Hom}(L,\mathbb{G}_m)\otimes\mathbb{Q}
  =\operatorname{Hom}(P,\mathbb{G}_m)\otimes\mathbb{Q}.
\]
称 $l\in\operatorname{Hom}(P,\mathbb{G}_m)^{*}$ \emph{严格支配}，是指对
$\alpha\in\Gamma\backslash\Delta$ 有 $l(p(\alpha))>0$。

对 $P$-丛 $\mathcal{F}$，记
$\deg\mathcal{F}\in\operatorname{Hom}(P,\mathbb{G}_m)^{*}$ 为如下泛函：
它把 $\varphi:P\to\mathbb{G}_m$ 对应到 $\mathcal{F}$ 经 $\varphi$
的前推的次数。称 $G$-丛为\emph{半稳定}的，是指对任何 $P\ne G$，
它都不来自严格支配次数的 $P$-丛。半稳定的 $G$-丛构成开子栈
$\operatorname{Bun}_G^{ss}\subset\operatorname{Bun}_G$。固定次数
$d\in\operatorname{Hom}(G,\mathbb{G}_m)$ 的半稳定 $G$-丛构成开子栈
$\operatorname{Bun}_G^{ss,d}\subset\operatorname{Bun}_G^{ss}$。若
$P\subset G$ 是含 $B$ 的抛物子群，
$d\in\operatorname{Hom}(P,\mathbb{G}_m)^{*}$ 严格支配，则记
$\operatorname{Shatz}_P^d$ 为次数为 $d$ 且相应 $L$-丛半稳定的
$P$-丛 $\mathcal{F}$ 之栈。已知自然态射
$\operatorname{Shatz}_P^d\to\operatorname{Bun}_G$ 是局部闭嵌入，且子栈%
\SourcePage{54}
$\operatorname{Shatz}_P^d$（$P,d$ 遍取）构成 $\operatorname{Bun}_G$
的一个分层，称为 \emph{Shatz 分层}。

记 $\operatorname{Nilp}_P^d(G)$（resp. $\operatorname{Nilp}^{ss}(G)$、
$\operatorname{Nilp}^{ss,d}(G)$）为
$\operatorname{Nilp}=\operatorname{Nilp}(G)$ 与
$\operatorname{Shatz}_P^d$（resp. $\operatorname{Bun}_G^{ss}$、
$\operatorname{Bun}_G^{ss,d}$）在 $\operatorname{Bun}_G$ 上的纤维积。
要证 $\operatorname{Nilp}(G)$ 的维数纯为
$\dim\operatorname{Bun}_G=0$，只需对每个 $P,d$ 证明
$\operatorname{Nilp}_P^d(G)$ 的维数纯为 $0$。设 $L$ 为 $P$ 的极大
约化商，$\mathfrak{p}:=\operatorname{Lie}P$，
$\mathfrak{l}:=\operatorname{Lie}L$。若 $\mathcal{F}$ 是具有严格支配
次数的 $P$-丛，且相应的 $L$-丛 $\overline{\mathcal{F}}$ 半稳定，则
$H^0(X,\mathfrak{g}_{\mathcal{F}})
=H^0(X,\mathfrak{p}_{\mathcal{F}})$，于是有自然映射
$\eta\mapsto\overline{\eta}$，从
$H^0(X,\mathfrak{g}_{\mathcal{F}})$ 到
$H^0(X,\mathfrak{l}_{\mathcal{F}})$。定义
$\pi:\operatorname{Nilp}_P^d(G)\to\operatorname{Nilp}^{ss,d}(L)$ 为
$(\mathcal{F},\eta)\mapsto(\overline{\mathcal{F}},\overline{\eta})$，
其中 $\eta\in H^0(X,\mathfrak{g}_{\mathcal{F}}\otimes\omega_X)
=H^0(X,\mathfrak{g}_{\mathcal{F}})$（因 $g=1$，$\omega_X$ 平凡）。
再次利用 $g=1$ 可知 $\pi$ 光滑，其纤维是 $0$ 维栈。故只需证
$\operatorname{Nilp}^{ss}(L)$ 的维数纯为 $0$。

$\operatorname{Nilp}^{ss}(L)$ 的点是二元组：半稳定的 $L$-丛
$\mathcal{F}$ 与幂零元 $\eta\in H^0(X,\mathfrak{l}_{\mathcal{F}})$。
因 $\mathfrak{l}_{\mathcal{F}}$ 是半稳定向量丛，
$\operatorname{ad}_\eta:\mathfrak{l}_{\mathcal{F}}\to
\mathfrak{l}_{\mathcal{F}}$ 的秩恒定。故 $\eta(x)$ 的共轭类不依赖
$x\in X$。对幂零共轭类 $C\subset\mathfrak{l}$，记
$\operatorname{Nilp}_C^{ss}(L)$ 为 $\operatorname{Nilp}^{ss}(L)$ 中
参数化满足 $\eta(x)\in C$ 的二元组 $(\mathcal{F},\eta)$ 的局部闭
子栈。只需对每个 $C$ 证明 $\operatorname{Nilp}_C^{ss}(L)$ 的维数
纯为 $0$。设 $Z(A)\subset L$ 是某 $A\in C$ 的中心化子，
$\mathfrak{z}(A):=\operatorname{Lie}Z(A)$。若
$(\mathcal{F},\eta)\in\operatorname{Nilp}_C^{ss}(L)$，则
$\eta\in\Gamma(X,C_{\mathcal{F}})
=\Gamma(X,(G/Z(A))_{\mathcal{F}})$ 在 $\mathcal{F}$ 上定义
$Z(A)$-结构。于是得到开嵌入
$\operatorname{Nilp}_C^{ss}(L)
\hookrightarrow\operatorname{Bun}_{Z(A)}$。最后，
$\operatorname{Bun}_{Z(A)}$ 的维数纯为 $0$：对任何 $Z(A)$-丛
$\mathcal{E}$ 有 $\chi(\mathfrak{z}(A)_{\mathcal{E}})
=\deg\mathfrak{z}(A)_{\mathcal{E}}=0$（注意 $G/Z(A)=C$ 带 $G$-不变
辛结构，故 $Z(A)$ 的伴随表示行列式平凡，从而
$\mathfrak{z}(A)_{\mathcal{E}}$ 平凡）。\hfill$\Box$

\bdp{2.10.5} \emph{命题 2.1.2 的证明.} 须证 (4) 对
$\mathcal{Y}=\operatorname{Bun}_G$ 成立，即对一切 $n>0$ 有
$\operatorname{codim}\{\mathcal{F}\in\operatorname{Bun}_G\mid
\dim H^0(X,\mathfrak{g}_{\mathcal{F}})=n\}>n$。

\SourcePage{55}

这等价于证明
\begin{equation}\tag{39}
  \dim(A(G)\backslash A^{0}(G))<\dim\operatorname{Bun}_G,
\end{equation}
其中 $A(G)$ 是二元组 $(\mathcal{F},s)$ 之栈，$\mathcal{F}\in
\operatorname{Bun}_G$，$s\in H^0(X,\mathfrak{g}_{\mathcal{F}})$；而
$A^{0}(G)\subset A(G)$ 是由方程 $s=0$ 定义的闭子栈。

置 $C:=\operatorname{Spec}(\operatorname{Sym}\mathfrak{g}^{*})^{G}$。
这是 $\mathfrak{g}$ 关于 $G$ 的伴随作用的仿射概形商；事实上
$C=W\backslash\mathfrak{h}$，其中 $\mathfrak{h}$ 是 $\mathfrak{g}$ 的
固定嘉当子代数，$W$ 是 Weyl 群。态射 $\mathfrak{g}\to C$ 诱导映射
$H^0(X,\mathfrak{g}_{\mathcal{F}})\to\operatorname{Mor}(X,C)=C$。
于是有典范态射 $f:A(G)\to C=W\backslash\mathfrak{h}$。对
$h\in\mathfrak{h}$，置 $A_h(G)=f^{-1}(\overline{h})$，其中
$\overline{h}\in W\backslash\mathfrak{h}$ 是 $h$ 的像。置
$G^{h}:=\{g\in G\mid ghg^{-1}=h\}$，$\mathfrak{g}^{h}:=
\operatorname{Lie}G^{h}=\{a\in\mathfrak{g}\mid[a,h]=0\}$。记
$\mathfrak{z}_h$ 为 $\mathfrak{g}^{h}$ 的中心。由于
$h\in\mathfrak{z}_h$，且 $\mathfrak{g}$ 中形如 $\mathfrak{z}_h$ 的
子代数只有有限个，(39) 由不等式
$\dim(A_h(G)\backslash A^{0}(G))
<\dim\operatorname{Bun}_G-\dim\mathfrak{z}_h$ 推出。故只需证明
\begin{equation}\tag{40}
  \dim A_h(G)<\dim\operatorname{Bun}_G-\dim\mathfrak{z}_h
  \qquad\text{当 } h\ne 0
\end{equation}
\begin{equation}\tag{41}
  \dim(A_0(G)\backslash A^{0}(G))<\dim\operatorname{Bun}_G.
\end{equation}

记 $Z_h$ 为 $G^{h}$ 的中心。下面说明：把 (41) 中的 $G$ 换成
$G^{h}/Z_h$，(40) 即可推出。事实上，有自然同构
\[
  A_0(G^{h})\xrightarrow{\sim}A_h(G^{h})\xrightarrow{\sim}A_h(G),
\]
以及显然的态射 $\varphi:A_0(G^{h})\to A_0(G^{h}/Z_h)$。
$\varphi$ 的非空纤维同构于 $\operatorname{Bun}_{Z_h}$，故
$\dim A_h(G)\le\dim\operatorname{Bun}_{Z_h}
+\dim A_0(G^{h}/Z_h)$。因
$\dim\operatorname{Bun}_{Z_h}=(g-1)\cdot\dim\mathfrak{z}_h$，而
(41) 蕴含 $\dim A_0(G^{h}/Z_h)
=(g-1)\cdot\dim(\mathfrak{g}^{h}/\mathfrak{z}_h)$，故
\[
  \dim A_h(G)\le(g-1)\cdot\dim\mathfrak{g}^{h}
  =\dim\operatorname{Bun}_G
  -(g-1)\cdot\dim(\mathfrak{g}/\mathfrak{g}^{h})
  \le\dim\operatorname{Bun}_G-\dim(\mathfrak{g}/\mathfrak{g}^{h}).
\]
最后，若 $h\ne 0$，则 $\dim(\mathfrak{g}/\mathfrak{g}^{h})
\ge 2\cdot\dim\mathfrak{z}_h>\dim\mathfrak{z}_h$。

为证 (41)，我们将证明：若 $Y\subset A_0(G)$ 是局部闭既约不可约
子栈，则 $\dim Y\le\dim\operatorname{Bun}_G$，且仅当
$Y\subset A^{0}(G)$ 时 $\dim Y=\dim\operatorname{Bun}_G$。对
$\xi\in H^0(X,\omega_X)$，考虑态射
$m_\xi:A_0(G)\to\operatorname{Nilp}$，
$(\mathcal{F},s)\mapsto(\mathcal{F},s\xi)$，其中 $\mathcal{F}\in
\operatorname{Bun}_G$，$s\in H^0(X,\mathfrak{g}_{\mathcal{F}})$。
这些态射%
\SourcePage{56}
$m_\xi$ 定义 $m:A_0(G)\times H^0(X,\omega_X)\to\operatorname{Nilp}$。
$m$ 的像含于某个局部闭既约不可约子栈
$Z\subset\operatorname{Nilp}$ 中。若 $\xi\ne 0$，则 $m_\xi$ 诱导
嵌入 $Y\hookrightarrow Z_\xi$，其中 $Z_\xi$ 是 $Z$ 的闭子栈，由满足
下述条件的二元组 $(\mathcal{F},\eta)\in
H^0(X,\mathfrak{g}_{\mathcal{F}}\otimes\omega_X)$ 组成：$\eta$ 在
子概形 $D_\xi:=\{x\in X\mid\xi(x)=0\}$ 上的限制为零。故
$\dim Y\le\dim Z_\xi\le\dim Z\le\dim\operatorname{Nilp}
=\dim\operatorname{Bun}_G$。若
$\dim Y=\dim\operatorname{Bun}_G$，则对所有非零
$\xi\in H^0(X,\omega_X)$ 有 $Z_\xi=Z$。这意味着所有
$(\mathcal{F},\eta)\in Z$ 都有 $\eta=0$，从而所有
$(\mathcal{F},s)\in Y$ 都有 $s=0$，即 $Y\subset A^{0}(G)$。
\hfill$\Box$

\bdsec{2.11. 论局部系统之栈} 记 $\mathcal{LS}_G$ 为 $X$ 上
$G$-局部系统之栈（$G$-局部系统即带联络的 $G$-丛）。Kapranov
[Kap97] 说明 $\mathcal{LS}_G$ 有导出版本 $R\mathcal{LS}_G$，后者是
DG 栈。利用 2.10 的结果，我们将证明：若 $g>1$ 且 $G$ 半单，则
$R\mathcal{LS}_G=\mathcal{LS}_G$。我们还将描述 $\mathcal{LS}_G$ 的
不可约分支之集。读者可跳过本节；其余部分不采用本节的结果。

\bdp{2.11.1} 固定 $x\in X$。记 $\mathcal{LS}_G^x$ 为 $X$ 上带联络
$\nabla$ 的 $G$-丛 $\mathcal{F}$ 之栈，该联络在 $x$ 有单极点。
% [Tutor Note: 原书 p.56 此处印作「G-biundles」，疑为「G-bundles」之笔误，底本如此，保留原样。]
记 $\mathcal{E}$ 为万有 $G$-丛（它定义在
$\mathcal{LS}_G^x\times X$ 上）在
$\mathcal{LS}_G^x=\mathcal{LS}_G^x\times\{x\}$ 上的限制。
% [Tutor Note: 原书 p.56 此处等式左端多写 $\mathcal{LS}_G^x$（应为 $\mathcal{LS}_G^x\times\{x\}$），底本如此，保留原样。]
$\nabla$ 在 $x$ 处的留数是截面
$R\in\Gamma(\mathcal{LS}_G^x,\mathfrak{g}_{\mathcal{E}})$，而
$\mathcal{LS}_G$ 是 $\mathcal{LS}_G^x$ 中由方程 $R=0$ 定义的闭子栈。
考虑开子栈 $\widetilde{\mathcal{LS}}{}_G^x\subset\mathcal{LS}_G^x$，
它参数化使 $\nabla:H^1(X,\mathfrak{g}_{\mathcal{F}})\to
H^1(X,\mathfrak{g}_{\mathcal{F}}\otimes\omega_X(x))$ 满射的二元组
$(\mathcal{F},\nabla)$。易见 $\widetilde{\mathcal{LS}}{}_G^x$ 是
纯维数 $(2g-1)\cdot\dim G$ 的光滑栈，且
$\mathcal{LS}_G\subset\widetilde{\mathcal{LS}}{}_G^x$。

把 $\mathfrak{g}_{\mathcal{E}}$ 看作 $\mathcal{LS}_G^x$ 上的栈。截面
$R,0\in\Gamma(\mathcal{LS}_G^x,\mathfrak{g}_{\mathcal{E}})$ 定义
$\mathfrak{g}_{\mathcal{E}}$ 的两个闭子栈；$R\mathcal{LS}$ 是它们在
导出意义下的交，而 $\mathcal{LS}_G$ 是通常的交。故下列条件等价：
\begin{enumerate}[label=\arabic*),leftmargin=2.2em]
\item $R\mathcal{LS}_G=\mathcal{LS}_G$；
\SourcePage{57}
\item $\mathcal{LS}_G$ 是纯维数 $(2g-2)\cdot\dim G$ 的局部完全交；
\item $\dim\mathcal{LS}_G\le(2g-2)\cdot\dim G$。
\end{enumerate}

下述命题表明：若 $g>1$ 且 $G$ 半单，则这些条件都满足。

\bdp{2.11.2} \emph{命题.} 设 $g>1$，$G$ 约化。则 $\mathcal{LS}_G$
是纯维数 $(2g-2)\cdot\dim G+l$ 的局部完全交，其中 $l$ 是 $G$ 的
中心的维数。

\emph{证明.} 设 $R$ 的含义与 2.11.1 中相同。显然
$R\in\Gamma(\mathcal{LS}_G^x,[\mathfrak{g},\mathfrak{g}]_{\mathcal{E}})$，
故只需证明
\begin{equation}\tag{42}
  \dim\mathcal{LS}_G\le(2g-2)\cdot\dim G+l.
\end{equation}
记 $G_{ad}$ 为 $G$ 模去中心所得的商。考虑投影
$p:\mathcal{LS}_G\to\operatorname{Bun}_{G_{ad}}$。若 $p$ 在
$G_{ad}$-丛 $\mathcal{F}$ 上的纤维非空，则其维数等于
$\dim T^{*}_{\mathcal{F}}\operatorname{Bun}_{G_{ad}}+l(2g-1)$，故
$\dim\mathcal{LS}_G\le\dim T^{*}\operatorname{Bun}_{G_{ad}}+l(2g-1)$。
最后，$\dim T^{*}\operatorname{Bun}_{G_{ad}}
\le\dim G_{ad}\cdot(2g-2)$，因为 $\operatorname{Bun}_{G_{ad}}$ 在
1.1.1 的意义下是良的（2.10.1 中已证）。\hfill$\Box$

\bdp{2.11.3} 记 $\operatorname{Bun}'_G\subset\operatorname{Bun}_G$
为 $\operatorname{Bun}_{G/[G,G]}$ 中含平凡丛的连通分支的原像。
$\mathcal{LS}_G\to\operatorname{Bun}_G$ 的像含于
$\operatorname{Bun}'_G$。

\bdp{2.11.4} \emph{命题.} 设 $g>1$，$G$ 约化。则
$\operatorname{Bun}'_G$ 的每个连通分支在 $\mathcal{LS}_G$ 中的原像
都非空且不可约。

于是 $\mathcal{LS}_G$ 的不可约分支由
\[
  \operatorname{Ker}(\pi_1(G)\to\pi_1(G/[G,G]))=\pi_1([G,G])
\]
参数化。

\emph{证明.} 考虑开子栈
$\operatorname{Bun}^0_{G_{ad}}\subset\operatorname{Bun}_{G_{ad}}$，
它参数化满足 $H^0(X,(\mathfrak{g}_{ad})_{\mathcal{F}})=0$ 的
$G_{ad}$-丛 $\mathcal{F}$（这是 $\operatorname{Bun}_{G_{ad}}$ 中
最大的 Deligne--Mumford 子栈）。记 $\operatorname{Bun}^{0\prime}_G$
为 $\operatorname{Bun}^0_{G_{ad}}$ 在 $\operatorname{Bun}'_G$ 中的
原像。

\SourcePage{58}

记 $\mathcal{LS}^0_G$ 为 $\operatorname{Bun}^0_{G_{ad}}$ 在
$\mathcal{LS}_G$ 中的原像。2.10.5 中已证
$\operatorname{Bun}_{G_{ad}}$ 在 1.1.1 的意义下是极良的，故
$\dim(T^{*}\operatorname{Bun}_{G_{ad}}
\backslash T^{*}\operatorname{Bun}^0_{G_{ad}})
<\dim T^{*}\operatorname{Bun}_{G_{ad}}$。(42) 的证明中所用的论证
表明 $\dim(\mathcal{LS}_G\backslash\mathcal{LS}^0_G)
<(2g-2)\cdot\dim G+l$。由 2.11.2 可知 $\mathcal{LS}^0_G$ 在
$\mathcal{LS}_G$ 中稠密。故只需证明
$\operatorname{Bun}^{0\prime}_G$ 的每个连通分支在
$\mathcal{LS}^0_G$ 中的原像都非空且不可约。这显然：态射
$\mathcal{LS}^0_G\to\operatorname{Bun}^{0\prime}_G$ 是
$T^{*}\operatorname{Bun}^{0\prime}_G$ 上的
torsor\footnote{torsor 结构依赖 $\mathfrak{g}$ 上一个不变数量积的
选取。}。\hfill$\Box$

\bdsec{2.12. 论 Beauville--Laszlo 定理} 本节实际上是一个附录，其中
解释 [BLa95] 主定理的整体化版本。2.3.7 用到这个版本，但并非本质地
用到。故读者可跳过本节。

\bdp{2.12.1} \emph{定理.} 设 $p:\widetilde{S}\to S$ 是概形态射，
$D\subset S$ 是有效 Cartier 除子。设
$\widetilde{D}:=p^{-1}(D)$ 是 $\widetilde{S}$ 中的 Cartier 除子，
且态射 $\widetilde{D}\to D$ 是同构。置 $U:=S\backslash D$，
$\widetilde{U}:=\widetilde{S}\backslash\widetilde{D}$。

记 $C$ 为由这样的拟凝聚 $\mathcal{O}_S$-模构成的范畴：它没有支撑在
$D$ 上的非零局部截面。记 $\widetilde{C}$ 为
$(\widetilde{S},\widetilde{D})$ 的类似范畴。记 $C'$ 为三元组
$(\mathcal{M}_1,\mathcal{M}_2,\varphi)$ 的范畴，其中
$\mathcal{M}_1$ 是拟凝聚 $\mathcal{O}_U$-模，
$\mathcal{M}_2\in\widetilde{C}$，$\varphi$ 是 $\mathcal{M}_1$ 与
$\mathcal{M}_2$ 到 $\widetilde{U}$ 的拉回之间的同构。
\begin{enumerate}[label=\arabic*),leftmargin=2.2em]
\item $p^{*}$ 把 $C$ 送到 $\widetilde{C}$，故有函子
  $F:C\to C'$，它把 $\mathcal{M}\in C$ 送到
  $(\mathcal{M}|_U,p^{*}\mathcal{M},\varphi)$，其中 $\varphi$ 是
  $\mathcal{M}|_U$ 与 $p^{*}\mathcal{M}$ 到 $\widetilde{U}$ 的拉回
  之间的自然同构。
\item $F:C\to C'$ 是范畴等价。
\item $\mathcal{M}\in C$ 局部有限型（resp. 平坦，resp. 有限秩局部
  自由），当且仅当 $\mathcal{M}|_U$ 与 $f^{*}\mathcal{M}$ 有此性质。
% [Tutor Note: 原书 p.58 此处印作 $f^{*}\mathcal{M}$；本节只定义了 $p$，疑为 $p^{*}\mathcal{M}$ 之笔误，底本如此，保留原样。]
\end{enumerate}

\SourcePage{59}

这一定理容易化归为下述情形：$S$ 与 $\widetilde{S}$
仿射\footnote{对任意 $x\in S$，存在 $x$ 的仿射邻域 $U$ 与开仿射
$\widetilde{U}\subset\widetilde{S}$，使得
$\widetilde{U}\subset p^{-1}(U)$，且 $\widetilde{U}\cap\widetilde{D}
=p^{-1}(U)\cap\widetilde{D}$。事实上，可设 $S$ 仿射且 $x\in D$。
设 $\widetilde{U}_1\subset\widetilde{S}$ 是 $x$ 在
$\widetilde{D}$ 中的原像的仿射邻域。则
$p(\widetilde{U}_1\cap\widetilde{D})$ 是 $x$ 在 $D$ 中的仿射邻域，
故对某个满足 $x\in U$ 的开仿射 $U\subset S$，它含 $U\cap D$。于是
$\widetilde{U}:=\widetilde{U}_1\times_S U$ 具有所需性质。}，且 $D$ 由
单个方程整体地定义（即 $S=\operatorname{Spec}A$，
$\widetilde{S}=\operatorname{Spec}\widetilde{A}$，
$D=\operatorname{Spec}A/fA$，$f\in A$ 不是零因子）。这种情形与
[BLa95] 中的处理相同（[BLa95] 假设
$\widetilde{A}=\widehat{A}:=A$ 关于 $f$-进拓扑的完备化，但
[BLa95] 只用到 $\widehat{A}$ 的两条性质：$f:\widehat{A}\to
\widehat{A}$ 的单射性，以及 $A/fA\to\widehat{A}/f\widehat{A}$ 的
双射性）。

\bdp{2.12.2} 设 $D$ 是概形 $S$ 的闭仿射子概形。记 $\widehat{S}$ 为
$S$ 沿 $D$ 的完备化，$\widehat{S}'$ 为 $\widehat{S}$ 上正则函数环的
谱（故 $\widehat{S}$ 是仿射形式概形，$\widehat{S}'$ 是相应的真
概形）。有态射 $\pi:\widehat{S}\to S$ 与
$i:\widehat{S}\to\widehat{S}'$。

\bdp{2.12.3} \emph{命题.} 至多存在一个态射
$p:\widehat{S}'\to S$ 使 $pi=\pi$。

\emph{证明.} 设对某 $p_1,p_2:\widehat{S}'\to S$ 有
$\pi=p_1i=p_2i$。设 $Y\subset\widehat{S}'$ 是对角
$\Delta\subset S\times S$ 在
$(p_1,p_2):\widehat{S}'\to S\times S$ 下的原像。则 $Y$ 是
$\widehat{S}'$ 的局部闭子概形，且对每个 $n$ 都含
$D\subset\widehat{S}'$ 的第 $n$ 个无穷小邻域。故
$(\overline{Y}\backslash Y)\cap D=\varnothing$，从而
$\overline{Y}\backslash Y=\varnothing$，即 $Y$ 闭。含 $D$ 的所有
无穷小邻域的 $\widehat{S}'$ 闭子概形等于 $\widehat{S}'$。故
$Y=\widehat{S}'$，$p_1=p_2$。\hfill$\Box$

\bdp{2.12.4} 设处于 2.12.2 的情形，且 $D\subset S$ 是有效 Cartier
除子。若存在 $p:\widehat{S}'\to S$ 使 $pi=\pi$，则
$p^{-1}(D)\subset\widehat{S}'$ 是 Cartier 除子，且态射
$p^{-1}(D)\to D$ 是同构。于是定理 2.12.1 适用。

\SourcePage{60}

\bdp{2.12.5} 设处于 2.12.2 的情形，且 $S$ 拟分离。则存在
$p:\widehat{S}'\to S$ 使 $pi=\pi$。我们所知的证明相当长：先处理
诺特情形，再用下述事实（Deligne 私下告知）：对任何拟紧拟分离概形
$S$，存在从 $S$ 到某个 $\mathbb{Z}$ 上有限型概形的仿射态射。

在 2.3.7 中，我们对 $S=X\otimes R$ 用到 $p:\widehat{S}'\to S$ 的
存在性，其中 $X$ 是我们的曲线，$R$ 是 $\mathbb{C}$-代数。故下述
结果够用。

\bdp{2.12.6} \emph{命题.} 设在 2.12.2 的情形中，$S$ 是
$\mathbb{P}^{n}\otimes R$ 的局部闭子概形（$R$ 为某环）。则存在
$p:\widehat{S}'\to S$ 使 $pi=\pi$。

\emph{证明.} 我们用 Jouanolou 技巧。设 $\mathbb{P}^{*}$ 是与
$\mathbb{P}=\mathbb{P}^{n}$ 对偶的射影空间，
$Z\subset\mathbb{P}\times\mathbb{P}^{*}$ 为关联对应，
$U:=(\mathbb{P}\times\mathbb{P}^{*})\backslash Z$。因态射
$U\to\mathbb{P}$ 是 $\mathbb{P}$ 上某向量丛上的 torsor，而
$\widehat{S}$ 是仿射形式概形，态射 $\widehat{S}\to\mathbb{P}$ 提升
为态射 $\widehat{S}\to U$。因 $U$ 仿射，
$\operatorname{Mor}(\widehat{S},U)=\operatorname{Mor}(\widehat{S}',U)$，
故得态射 $\widehat{S}'\to U$。合成
$\widehat{S}'\to U\to\mathbb{P}$ 给出态射
$f:\widehat{S}'\to\mathbb{P}\otimes R$。局部闭子概形
$f^{-1}(S)\subset\widehat{S}'$ 对每个 $n$ 都含
$D\subset\widehat{S}'$ 的第 $n$ 个无穷小邻域，故
$f^{-1}(S)=\widehat{S}'$（参看 2.12.3），且 $f$ 诱导态射
$p:\widehat{S}'\to S\subset\mathbb{P}\otimes R$。显然
$pi=\pi$。\hfill$\Box$

\emph{注.} 也可用下述方式证明本命题：把态射
$\widehat{S}\to\mathbb{P}^{n}$ 解释为二元组
$(\mathcal{M},\varphi)$，其中 $\mathcal{M}$ 是 $\widehat{S}$ 上的
可逆层，$\varphi$ 是满态射 $\mathcal{O}^{n+1}\to\mathcal{M}$。然后
证明 $(\mathcal{M},\varphi)$ 延拓为 $\widehat{S}'$ 上的二元组
$(\mathcal{M}',\varphi')$。当然，此证明与基于 Jouanolou 技巧的证明
本质等价。
